\acp{SNN} represent an alternative type of neural networks that model neuronal behavior more closely to biological systems by processing time-dependent, spike-based signals.
In contrast, traditional \acp{ANN} operate on real-valued, continuous signals and follow a synchronous computation model.
This paper compares the architectural principles of \acp{SNN} and \acp{ANN}, focusing on neuron models, encoding strategies, signal propagation and activation mechanisms and learning paradigms.
It further investigates the suitability of \acp{GPU} for executing \acp{SNN}, despite their inherent mismatch with the event-driven and sparse activity of \acp{SNN}.
Several optimization techniques and \ac{GPU}-based frameworks, such as Temporal Fusion, Brian2CUDA, and Norse are presented based on findings from the literature.
The results suggest that while neuromorphic hardware remains the most suitable execution platform for \acp{SNN}, modern \acp{GPU} methodologies can serve as a viable alternative for both simulation and training.